% Generated draft: requires mathematical and source review.
\documentclass[11pt]{article}
\usepackage[margin=25mm]{geometry}
\usepackage{fontspec}
\setmainfont{Latin Modern Roman}
\usepackage{amsmath,amssymb,mathtools,unicode-math}
\setmathfont{Latin Modern Math}
\usepackage{enumitem,needspace,xurl,hyperref}
\setlength{\parindent}{0pt}
\setlength{\parskip}{5pt}
\setlength{\emergencystretch}{3em}
\newcommand{\N}{\mathbb{N}}
\newcommand{\Z}{\mathbb{Z}}
\newcommand{\Q}{\mathbb{Q}}
\newcommand{\R}{\mathbb{R}}
\newcommand{\C}{\mathbb{C}}
\newcommand{\F}{\mathbb{F}}
\newcommand{\A}{\mathbb{A}}
\newcommand{\PP}{\mathbb{P}}
\newcommand{\E}{\mathbb{E}}
\newcommand{\eps}{\varepsilon}
\newcommand{\dd}{\,\mathrm{d}}
\begin{document}
\section*{UnsolvedMath ID: 30002467}
\textbf{Source problem code:} catalogue entry
\subsection{Definitions and mathematical statement}
For integers \(k\ge 2\) and \(n\ge 0\), a simple \(k\)-uniform hypergraph \(H\) consists of a vertex set \(V(H)\) and an edge set \(E(H)\subseteq\binom{V(H)}{k}\). For a fixed \(k\)-uniform hypergraph \(F\), define
\[
\operatorname{ex}(n;F)=\max\bigl\{|E(H)|:\ |V(H)|=n,\ H\text{ contains no subhypergraph isomorphic to }F\bigr\}.
\]
Containment here is not required to be induced. These are the usual simple-hypergraph conventions; the supplied record does not explicitly specify them or the parameter ranges.

For \(m\ge k\), define the tight path \(T_m^k\) on distinct vertices \(v_1,\ldots,v_m\) by
\[
E(T_m^k)=\bigl\{\{v_i,v_{i+1},\ldots,v_{i+k-1}\}:1\le i\le m-k+1\bigr\}.
\]
Thus it has \(m\) vertices and \(m-k+1\) edges. For \(m\ge k+1\), define the tight cycle \(D_m^k\) on distinct vertices \(v_0,\ldots,v_{m-1}\) by
\[
E(D_m^k)=\bigl\{\{v_i,v_{i+1},\ldots,v_{i+k-1}\}:i\in\Z/m\Z\bigr\},
\]
where subscripts are read modulo \(m\). In this range it has \(m\) vertices and \(m\) edges.

\textbf{Source question.} Determine
\[
\operatorname{ex}(n;P_\ell^k)\qquad\text{and}\qquad\operatorname{ex}(n;C_\ell^k)
\quad\text{when }k\mid\ell,
\]
where the source calls \(P_\ell^k\) a tight path and \(C_\ell^k\) a tight cycle.

\textbf{Unresolved indexing and scope.} The record does not say whether \(\ell\) counts vertices or edges of the path. Under vertex indexing, the path target is \(T_\ell^k\), with \(\ell\ge k\); under edge indexing it is \(T_{\ell+k-1}^k\), with \(\ell\ge1\). These give different questions even with \(k\mid\ell\). For the nondegenerate cycle definition above, either count gives \(C_\ell^k=D_\ell^k\) for \(\ell\ge k+1\). The source does not specify how to treat smaller or degenerate cycles, including \(\ell=k\), or whether \(k=1\) is intended. Consequently, the displayed definitions clarify conventional interpretations without resolving the source's missing conventions. The request is to determine these extremal functions for the intended admissible parameters; the source gives no restriction to sufficiently large \(n\) and does not specify whether an exact or asymptotic determination is sought.
\subsection{Short English statement}
Determine the maximum number of edges in an \(n\)-vertex \(k\)-uniform hypergraph avoiding a specified tight path or tight cycle when \(k\mid\ell\). The path length convention and the intended parameter boundaries require clarification. The supplied record labels the general questions open; this draft does not independently verify that status.
\subsection{Sources}
\begin{itemize}
\item Local UnsolvedMath registry, permanent ID \(30002467\), source problem code \texttt{catalogue entry}, title \emph{Extremal Numbers of Tight Hypergraph Paths and Cycles}. The record belongs to the Oberwolfach Reports Open Problems collection and assigns the year 2014. No individual proposer is supplied.
\item Supplied primary report citation: \emph{Combinatorics} (2014), Oberwolfach Reports, pp.~5--90, DOI 10.4171/owr/2014/01. Also listed as ``Original OWR record'': \url{https://doi.org/10.4171/owr/2014/01}. The supplied background describes these reports as documenting mathematical workshops and open-problem discussions; no precise location within the report is supplied.
\item Supplied literature reference: \emph{Tur\'an problems for paths and cycles}, \url{https://www.sciencedirect.com/science/article/pii/S0097316514001204}.
\item Supplied literature reference: \emph{Linear hypergraph cycles}, \url{https://arxiv.org/abs/1302.2387}.
\item Literature assessment checked at \texttt{2026-08-21}: the supplied curation labels the problem open and states that primary work settles important fixed-uniformity path and cycle families while the general extremal-number questions remain unresolved. This is supplied dated curation derived from the OWR record, not an independently verified literature assessment. Its time-sensitive status and cited claims require independent verification; the listed works' applicability to the tight-hypergraph formulations also requires review.
\end{itemize}
\end{document}
